% L3-1 (2026-10-28): paste into §6.5 (numerical_results.tex), before \subsubsection{Point resonator on the plate}.
% Figure: plots/L3-1/oblique_plate.png -> figures/L3_1_oblique_plate.png (scripts/export_paper_assets.jl).
% Data: data/exp_pro/L3-1/{results,ladder,checks}.csv (144/144 checks).

\subsubsection{Bare plate under oblique incidence}\label{subsubsec:num_L3_1}

\paragraph{Test description.} The plate covers the whole surface of a box of depth $H=10~\si{\meter}$, with the parameters of Table~\ref{tab:dispersion_test_parameters} and $\nu=0$. The box is periodic in $y$ with width $W=\lambda_p/\sin\theta$. A line load $F\mathrm{e}^{\mathrm{i}k_yy}$ with $k_y=k_0\sin\theta$ acts on a strip of two cells across the plate, and viscoelastic sponges absorb the waves at both ends of the box. Since $\mathcal{D}_0$ depends on $|\boldsymbol{k}|$ only, the load radiates two oblique plate waves with $k_x=(k_0^2-k_y^2)^{1/2}$. Closing the $k_x$ integral on its poles gives
\begin{equation}
    \eta(x,y)=F\mathrm{e}^{\mathrm{i}k_yy}\,\mathrm{i}\sum_m\frac{k_m}{k_{x,m}}\frac{\mathrm{e}^{\mathrm{i}k_{x,m}|x|}}{\mathcal{D}_0'(k_m)},\qquad k_{x,m}=(k_m^2-k_y^2)^{1/2},
\end{equation}
so the radiated amplitude is $\gamma_\theta=\gamma_r/\cos\theta$. Each wave carries the flux $\tfrac14\omega\mathcal{D}_0'(k_0)|A|^2\cos\theta$ in $x$, and the input power $\tfrac12\omega\,\mathrm{Im}(\bar F\eta)$ equals the net outgoing flux. With $k_y\neq0$ the fluid evanescent modes decay at least as fast as $(\kappa_1^2+k_y^2)^{1/2}$. The fields are $Q_2$, with $h=\lambda_p/8$ horizontally and $8$ graded cells vertically (top cell $0.2/k_0$). The channel holds sponges of $4\lambda_p$, gaps of $5/\kappa_1$ and $3\lambda_p$ on each side of the load: at $\omega=10~\si{\radian\per\second}$ and $\theta=30^\circ$ this gives $N=3.8\times10^5$. We run $\theta\in\{0^\circ,30^\circ,60^\circ\}$ at $\omega=10~\si{\radian\per\second}$, and $\theta\in\{0^\circ,30^\circ\}$ at $\omega=5$ and $2~\si{\radian\per\second}$. The wavenumber and the amplitudes come from a fit of the $\mathrm{e}^{\mathrm{i}k_yy}$ component at the cell vertices, where a discrete plane wave is exactly exponential. This mesh cannot be refined within the memory available, so the mesh study uses a shorter channel (sponges $2\lambda_p$, gaps $5/(\kappa_1^2+k_y^2)^{1/2}$) with $h=\lambda_p/6$, $\lambda_p/8$ and $\lambda_p/12$. The vertical refinement is that of the vertical section (Section~\ref{subsubsec:num_L1_1}).

\paragraph{Results.} All 144 checks pass. The residuals are below $5\times10^{-9}$. Two point loads are reciprocal to $5\times10^{-10}$. The field projected onto $\mathrm{e}^{\mathrm{i}k_yy}$ fits two discrete plane waves and the evanescent modes to $6\times10^{-9}$ at the cell vertices ($2\times10^{-6}$ at normal incidence). On the reference mesh $h=\lambda_p/8$, $k_h=(k_{x,h}^2+k_y^2)^{1/2}$ is within $0.42\%$ of $k_0$ at $\theta=30^\circ$ and within $0.45\%$ at $60^\circ$, against $0.93\%$ at normal incidence ($\omega=10~\si{\radian\per\second}$). At $\omega=5$ and $2~\si{\radian\per\second}$ the errors are $0.32\%$ and $0.02\%$ at $30^\circ$. The radiated amplitude follows $\gamma_r/\cos\theta$, and the net flux balances the input power, within $2.3\%$ at $30^\circ$ and $3.8\%$ at $0^\circ$ (Figure~\ref{fig:num_L3_1}c). The shorter channel gives the same $k_h$ to $10^{-8}$. The error is therefore the same discretisation error as in the vertical section, and it is smaller for oblique waves. In the mesh study (Figure~\ref{fig:num_L3_1}b) this anisotropy, $|k_h(30^\circ)/k_h(0^\circ)-1|$, falls from $0.51\%$ to $0.23\%$ between $h=\lambda_p/8$ and $\lambda_p/12$, a rate of $1.9$. The errors in $k_h$, in the amplitude and in the energy balance converge at the same rate for both angles, and the rate stays at $1.9$ up to $h=\lambda_p/16$ at normal incidence. Within the discretisation error, the three-dimensional formulation thus reproduces the isotropic dispersion relation of the host and the oblique Green function.

\begin{figure}[t]
    \centering
    \includegraphics[width=\linewidth]{figures/L3_1_oblique_plate.png}
    \caption{Bare plate under oblique incidence, periodic in $y$. (a) $\eta$ along $y=W/3$, normalised by $F\gamma_\theta$, for $\theta=30^\circ$, $\omega=10~\si{\radian\per\second}$, $h=\lambda_p/8$; shaded: fit windows; dashed: start of the sponges. (b) Mesh study in the short channel: $|k_h/k_0-1|$ at $\theta=0^\circ$ and $30^\circ$ and the anisotropy $|k_h(30^\circ)/k_h(0^\circ)-1|$ against $h/\lambda_p$. (c) Radiated amplitude $\gamma_h/\gamma_\theta$ and net flux over input power against $\theta$, $h=\lambda_p/8$, for $\omega\in\{10,5,2\}~\si{\radian\per\second}$.}
    \label{fig:num_L3_1}
\end{figure}
